\section{Выводы}
Выполнив лабораторные работы по курсу \enquote{Информационный поиск}, я получил полное практическое понимание архитектуры поисковой системы — от сбора данных до пользовательского интерфейса. Реализация всех компонентов «с нуля» позволила глубоко разобраться в деталях, которые скрыты при использовании готовых решений, таких как Elasticsearch или Whoosh.
Во-первых, я освоил современные подходы к сбору данных: вместо медленного обхода веб-страниц я научился работать с готовыми дампами с Archive.org, эффективно распаковывая архивы \texttt{.7z} и извлекая структурированные данные из XML. Это показало, что для многих задач не нужно писать сложного краулера — достаточно использовать открытые источники с качественным контентом, такими как Stack Exchange.
Во-вторых, реализация токенизатора и стеммера на C++ без использования STL (кроме \texttt{std::string} и \texttt{std::vector}) позволила оценить важность контроля над производительностью и памятью. Я лучше понял, как устроены базовые операции: чтение файла в память, нормализация текста, разбиение на слова. Стеммер, вдохновлённый алгоритмом Портера, продемонстрировал, насколько даже простые правила могут эффективно нормализовать лексику и повысить точность поиска.
В-третьих, построение бинарных индексов (\texttt{inverted\_index.bin}, \texttt{forward\_index.bin}) показало, насколько важна компактность и скорость доступа. Хранение данных в сериализованном виде позволяет быстро загружать индексы при старте программы и минимизирует использование памяти. Сортировка терминов по алфавиту открывает возможность для будущего улучшения — например, бинарного поиска или использования trie.
Наконец, реализация булева поиска с поддержкой \texttt{\&\&}, \texttt{||}, \texttt{!} и скобок показала силу алгоритмов слияния отсортированных списков. Даже на корпусе из 30\,000 документов время ответа составляет доли миллисекунды. Интеграция с веб-интерфейсом через Flask превратила консольную утилиту в удобный сервис, доступный из любого браузера, с постраничной навигацией и чистым дизайном.
Полученный опыт имеет прямое применение в реальных проектах: от корпоративного поиска до анализа технической документации и создания ассистентов. Теперь я не просто использую поисковые технологии — я понимаю, как они устроены внутри, и могу адаптировать их под конкретные задачи.
\pagebreak
